\documentclass[10pt,a4paper]{amsart}
\usepackage[margin=19mm]{geometry}
\usepackage{amsmath,amssymb,mathtools}
\usepackage[scr=rsfs,cal=euler]{mathalfa}
\usepackage{xfrac}
\usepackage[unicode,hidelinks,bookmarks=false]{hyperref}
\newcommand{\ie}{i.e.\ }
\newcommand{\cf}{cf.\ }
\newcommand{\resp}{respectively\ }
\newcommand{\Subsection}{\subsection}
\providecommand{\nobreakdash}{\nobreak\hbox{-}}
\setlength{\emergencystretch}{2em}
\multlinegap0pt
\usepackage{bm}
\usepackage{bbm}
\usepackage{dsfont}
\usepackage{xspace}
\usepackage{cancel}
\usepackage{mathtools}
\usepackage{amsthm}
\makeatletter
\@ifundefined{lemma}{\newtheorem{lemma}{Lemma}}{}
\@ifundefined{theorem}{\newtheorem{theorem}{Theorem}}{}
\makeatother
\title[Linearly Stable Generalizations of ESFR Schemes]{Linearly Stable Generalizations of ESFR Schemes}
\author{Mathias Dufresne-Pich\'e, Siva Nadarajah}
\date{Source: arXiv:2510.15043v1 (CC BY 4.0)}
\begin{document}
\maketitle

\noindent\emph{Source and licence.} Linearly Stable Generalizations of ESFR Schemes. Version arXiv:2510.15043v1 (CC BY 4.0), distributed under the Creative Commons Attribution 4.0 International licence, as recorded in the arXiv licence metadata for this exact version and shown on the abstract page https://arxiv.org/abs/2510.15043v1. Full rights notice: licenses/arxiv-2510.15043-CC-BY-4.0.txt.\par\smallskip

\noindent\textit{Editorial scope.} Counted unit: the Theorem whose source label is \texttt{None} in the article (source0.tex lines 899--901, source label \texttt{None}), delivered together with the article's own complete original proof of it (source0.tex lines 902--956, one \texttt{proof} environment of 55 lines). No mathematics is written, repaired or re-derived: statement and proof are the article's own text line for line. Required context retained uncounted: lem:equivalence (lines 849--861); eq:equivalence1 (lines 891--895); eq:equivalence2 (lines 891--895); eq:equivalence3 (lines 891--895). Every definition, display and label of those lines is retained unchanged. Omitted from this package and not redistributed: 1--848, 862--890, 896--898, 957--993. Nothing inside a retained excerpt is omitted or altered: every retained body is delivered exactly as the article's lines are written, so no omission receipt of any kind accompanies this package. Citations: the retained excerpts carry no citation command at all, so no bibliography entry of the article is transcribed and no reference is redistributed. Reference adaptation: none was needed and none was made; every reference inside the delivered unit resolves inside it, so no unresolved reference or citation is produced. Printed slips kept literal: no printed word, symbol, hypothesis or number of the counted statement or of its proof was altered. Layout and class: the article's own class is \texttt{elsarticle} with packages including hyperref, graphicx, svg, caption, subcaption, amssymb, amsmath, amsthm, bm; the wrapper is the lane's pinned \texttt{amsart} editorial wrapper at 10pt on A4, which restates the theorem-like environments the retained text needs and adds the article's own macro definitions that the retained text needs (none), with extra wrapper packages (bm, bbm, dsfont, xspace, cancel, mathtools). The delivered numbering is the unit's own. Assets: no retained span contains a figure environment, a graphics-inclusion command, a PDF-page inclusion, a table or a TikZ picture; no image file is copied into this package, the package declares no asset dependency and no image file is redistributed with it. Licence: version arXiv:2510.15043v1 of this article is distributed under the Creative Commons Attribution 4.0 International licence, as recorded in the arXiv licence metadata for this exact version and shown on the abstract page https://arxiv.org/abs/2510.15043v1. A full rights notice is delivered at licenses/arxiv-2510.15043-CC-BY-4.0.txt. Characters: every retained excerpt is pure ASCII, so the delivered engine, XeLaTeX with its default Latin Modern text font, reports no missing character.
\begin{lemma}
\label{lem:equivalence}
    A FDG scheme can be expressed in FR form if and only if its associated filtering matrix $\mathbf{K} \in \mathbb{R}^{(p+1) \times (p+1)}$ is such that $\mathbf{M} + \mathbf{K}$ is invertible and
    \begin{equation}
        \mathbf{K}\mathbf{D} = 0.
    \end{equation}
    Additionally, a FR scheme can be expressed in FDG form if and only if its associated matrix $\mathbf{Q} \in \mathbb{R}^{(p+1) \times (p+1)}$ is such that $\mathbf{M} + \mathbf{Q}$ is invertible and $\exists \mathbf{K} \in \mathbb{R}^{(p+1) \times (p+1)}$ with $\mathbf{M} + \mathbf{K}$ invertible satisfying
    \begin{align}
    \mathbf{K}\mathbf{D} &= 0, \\
    (\mathbf{M} + \mathbf{K})^{-1} \mathbf{r} &= (\mathbf{M} + \mathbf{Q})^{-1} \mathbf{r},\\
    (\mathbf{M} + \mathbf{K})^{-1} \mathbf{l} &= (\mathbf{M} + \mathbf{Q})^{-1} \mathbf{l}.
\end{align}
\end{lemma}

\begin{align}
    \mathbf{K}\mathbf{D} &= 0 \label{eq:equivalence1}, \\
    (\mathbf{M} + \mathbf{K})^{-1} \mathbf{r} &= (\mathbf{M} + \mathbf{Q})^{-1} \mathbf{r}, \label{eq:equivalence2}\\
    (\mathbf{M} + \mathbf{K})^{-1} \mathbf{l} &= (\mathbf{M} + \mathbf{Q})^{-1} \mathbf{l} \label{eq:equivalence3}.
\end{align}

\begin{theorem}
    If a given numerical scheme that can be expressed in FR and FDG form is such that $\mathbf{K} = \mathbf{K}^T$ or $\mathbf{J}\mathbf{Q} = \mathbf{Q}\mathbf{J}$, then it must be an ESFR scheme, \textit{ie.}, ESFR schemes are the only numerical schemes satisfying the conditions from Lemma \ref{lem:equivalence} with symmetric filtering operators or symmetric correction functions.
\end{theorem}

\begin{proof}
Without loss of generality, in the following, we assume that all computations are performed in the Legendre basis. In this case, recall that $\mathbf{J}_{ij} = (-1)^{i+1} \delta_{ij}$. We first note that the only possible form for a matrix $\mathbf{K}$ satisfying Eq.(\ref{eq:equivalence1}) is
\begin{equation}
    \forall_{i=0}^p \forall_{j=0}^{p-1} : (\mathbf{K})_{ij} = 0.
    \label{eq:K_constraint}
\end{equation}
Let $\mathbf{K} = \mathbf{K}^T$. Then, the only non-zero entry in $\mathbf{K}$ is $(\mathbf{K})_{pp}$ and hence $\mathbf{K}$ is the usual ESFR filtering matrix. This completes the first part of the proof.

Using Eq.(\ref{eq:K_constraint}), direct computations shows that $(\mathbf{M}+\mathbf{K})^{-1}$ is given by
\begin{align}
    &(\mathbf{M} + \mathbf{K})^{-1} \nonumber
    =\\
    &\begin{bmatrix}
        (\mathbf{M}^{-1})_{00} &0 & & \cdots& -(\mathbf{M}^{-1})_{00}(\mathbf{M}^{-1})_{pp}(\mathbf{K})_{00}/(1+(\mathbf{M}^{-1})_{pp}(\mathbf{K})_{pp}) \\
        0&(\mathbf{M}^{-1})_{11}&0 & \cdots& -(\mathbf{M}^{-1})_{11}(\mathbf{M}^{-1})_{pp}(\mathbf{K})_{11}/(1+(\mathbf{M}^{-1})_{pp}(\mathbf{K})_{pp}) \\
        \vdots&\vdots&\vdots && \vdots \\
        0&0&0&\cdots&(\mathbf{M}^{-1})_{pp}/(1+(\mathbf{M}^{-1})_{pp} (\mathbf{K})_{pp})
    \end{bmatrix}
    \label{eq:MK_form}
\end{align}
Let us now assume that $\mathbf{J}\mathbf{Q} = \mathbf{Q}\mathbf{J}$. Then, knowing that $\mathbf{l} = -\mathbf{J}\mathbf{r}$ and that $\mathbf{J} = \mathbf{J}^{-1}$, we must have
\begin{align}
    (\mathbf{M} + \mathbf{Q})\mathbf{J} = \mathbf{J}(\mathbf{M}+\mathbf{Q})
    &\implies
    \mathbf{J}(\mathbf{M} + \mathbf{Q})^{-1} = (\mathbf{M} + \mathbf{Q})^{-1}\mathbf{J} \\
    &\implies
    \mathbf{J}(\mathbf{M} + \mathbf{Q})^{-1}\mathbf{r} = (\mathbf{M} + \mathbf{Q})^{-1}\mathbf{J}\mathbf{r} \\
    &\implies
    \mathbf{J}(\mathbf{M} + \mathbf{Q})^{-1}\mathbf{r} = -(\mathbf{M} + \mathbf{Q})^{-1}\mathbf{l}
\end{align}
Using Eq.(\ref{eq:equivalence2}) and Eq.(\ref{eq:equivalence3}), this implies that
\begin{equation}
    \mathbf{J}(\mathbf{M} + \mathbf{K})^{-1}\mathbf{r} = -(\mathbf{M} + \mathbf{K})^{-1}\mathbf{l}.
    \label{eq:intermediate}
\end{equation}
Substituting Eq.(\ref{eq:MK_form}) in Eq.(\ref{eq:intermediate}), one obtains
\begin{align}
    &\begin{bmatrix}
        (-1)^1(\mathbf{M}^{-1})_{00} +(-1)^1(\mathbf{M}^{-1})_{00}(\mathbf{M}^{-1})_{pp}(\mathbf{K})_{00}/(1+(\mathbf{M}^{-1})_{pp}(\mathbf{K})_{pp}) \\
        (-1)^2(\mathbf{M}^{-1})_{11} +(-1)^3(\mathbf{M}^{-1})_{11}(\mathbf{M}^{-1})_{pp}(\mathbf{K})_{11}/(1+(\mathbf{M}^{-1})_{pp}(\mathbf{K})_{pp}) \\
        (-1)^3(\mathbf{M}^{-1})_{11} +(-1)^p(\mathbf{M}^{-1})_{22}(\mathbf{M}^{-1})_{pp}(\mathbf{K})_{22}/(1+(\mathbf{M}^{-1})_{pp}(\mathbf{K})_{pp}) \\
        \vdots \\
        (-1)^{p+1}(\mathbf{M}^{-1})_{pp}/(1+(\mathbf{M}^{-1})_{pp} (\mathbf{K})_{pp})
    \end{bmatrix}
    =\nonumber \\
    &\begin{bmatrix}
        (-1)^1(\mathbf{M}^{-1})_{00} + (-1)^0(\mathbf{M}^{-1})_{00}(\mathbf{M}^{-1})_{pp}(\mathbf{K})_{00}/(1+(\mathbf{M}^{-1})_{pp}(\mathbf{K})_{pp}) \\
        (-1)^2(\mathbf{M}^{-1})_{11} +(-1)^1(\mathbf{M}^{-1})_{11}(\mathbf{M}^{-1})_{pp}(\mathbf{K})_{11}/(1+(\mathbf{M}^{-1})_{pp}(\mathbf{K})_{pp}) \\
        (-1)^3(\mathbf{M}^{-1})_{22} +(-1)^2(\mathbf{M}^{-1})_{22}(\mathbf{M}^{-1})_{pp}(\mathbf{K})_{11}/(1+(\mathbf{M}^{-1})_{pp}(\mathbf{K})_{pp}) \\
        \vdots \\
        (-1)^{p+1}(\mathbf{M}^{-1})_{pp}/(1+(\mathbf{M}^{-1})_{pp} (\mathbf{K})_{pp})
    \end{bmatrix}
\end{align}
Clearly, this constraint can only be satisfied if $(\mathbf{K})_{pp}$ is the only non-zero entry of $\mathbf{K}$. Hence $\mathbf{K}$ is the ESFR filetring matrix, which completes the second part of the proof.
\end{proof}

\end{document}
